\clearpage
\chapter*{GLOSARIO}
\addcontentsline{toc}{chapter}{GLOSARIO}

\noindent Los términos legales se definen conforme a la LOPDP y a su Reglamento General (Tabla~\ref{tab:terminos_legales}); aquí se reúnen, además, los términos técnicos y las siglas empleados en el documento.

\begin{description}
  \item[AMIE] Archivo Maestro de Instituciones Educativas del Ministerio de Educación.
  \item[CES] Consejo de Educación Superior.
  \item[CSIRT] Equipo de respuesta a incidentes de seguridad informática; desde 2026, destinatario informativo de las notificaciones de vulneraciones (LOPDP, art.~43).
  \item[Delphi] Técnica de consulta a expertos en rondas anónimas con retroalimentación controlada, hasta alcanzar consenso.
  \item[DPD] Delegado de protección de datos personales (LOPDP, arts.~48 y 49).
  \item[DSR] Investigación en ciencia del diseño (\textit{Design Science Research}).
  \item[EIPD] Evaluación de impacto relativa a la protección de datos (LOPDP, art.~42).
  \item[Escenario de riesgo] Combinación de una comunidad de amenaza, un vector, una o más vulnerabilidades y el impacto resultante sobre el titular y sobre el responsable.
  \item[IES] Institución de educación superior.
  \item[LOFC] Ley Orgánica para el Fortalecimiento de la Ciberseguridad (2026).
  \item[LOPDP] Ley Orgánica de Protección de Datos Personales (2021).
  \item[Mapa de calor] Representación gráfica del número de escenarios por nivel de probabilidad e impacto; se usa para comunicar, no para calibrar.
  \item[Puntaje de Brier] Medida de la calibración de un estimador probabilístico: media de los cuadrados de la diferencia entre la probabilidad asignada y el resultado observado.
  \item[RAT] Registro de actividades de tratamiento (RLOPDP, arts.~38 y 39).
  \item[\textit{Rationale}] Justificación documentada de un valor de probabilidad o de impacto, exigida por la guía de la SPDP.
  \item[RIC] Rango intercuartílico: diferencia entre el tercer y el primer cuartil de las respuestas de un panel.
  \item[Riesgo inherente] Nivel de riesgo antes de aplicar los controles propuestos.
  \item[Riesgo residual] Nivel de riesgo estimado tras aplicar los controles propuestos.
  \item[RLOPDP] Reglamento General a la Ley Orgánica de Protección de Datos Personales (Decreto Ejecutivo N.º 904, 2023).
  \item[SISPDP] Sistema Nacional de Registro de Datos Personales, canal digital de notificación de vulneraciones a la SPDP.
  \item[SPDP] Superintendencia de Protección de Datos Personales.
  \item[Término] Plazo legal expresado en días que la LOPDP fija para la notificación de vulneraciones (arts.~43 y 46).
  \item[UTI] Universidad Tecnológica Indoamérica.
  \item[V de Aiken] Coeficiente de validez de contenido que resume el acuerdo de un grupo de jueces sobre un enunciado.
\end{description}
